% ==============================================================================
% Chapitre 4 : Tests et Validation - Rapport PFE - ISET Nabeul
% ==============================================================================

\chapter{Tests et Validation}
\label{chap:tests-validation}

\vspace{0.5cm}

\section*{Introduction}

Apr\`{e}s avoir pr\'{e}sent\'{e} la r\'{e}alisation technique de notre robot de tests
automatis\'{e}s dans le chapitre pr\'{e}c\'{e}dent, ce chapitre est consacr\'{e} \`{a}
l'ex\'{e}cution des diff\'{e}rentes suites de tests et \`{a} l'analyse de leurs
r\'{e}sultats. Nous d\'{e}taillons ici l'environnement de test utilis\'{e}, les
r\'{e}sultats obtenus pour chaque niveau de la pyramide de tests, ainsi que
les probl\`{e}mes rencontr\'{e}s et les solutions apport\'{e}es. Ce chapitre
constitue la preuve concr\`{e}te de la valeur ajout\'{e}e de notre travail et
valide l'ensemble de la d\'{e}marche adopt\'{e}e.

% ==============================================================================
\section{Environnement de test}
\label{sec:env-test}
% ==============================================================================

La mise en place d'un environnement de test fiable et reproductible est
un pr\'{e}requis essentiel pour garantir la coh\'{e}rence des r\'{e}sultats. Nous
avons utilis\'{e} deux environnements compl\'{e}mentaires~: un environnement
local pour le d\'{e}veloppement et le d\'{e}bogage, et un environnement
d'int\'{e}gration continue pour l'ex\'{e}cution automatis\'{e}e.

\subsection{Environnement local}

L'environnement de d\'{e}veloppement local a \'{e}t\'{e} configur\'{e} sur un poste
\'{e}quip\'{e} de macOS avec les composants suivants~:

\begin{table}[H]
\centering
\caption{Configuration de l'environnement local de test}
\label{tab:env-local}
\begin{tabularx}{\textwidth}{|l|X|}
\hline
\rowcolor{speedlineBlue!10}
\textbf{Composant} & \textbf{Version / D\'{e}tails} \\
\hline
Syst\`{e}me d'exploitation & macOS (Apple Silicon / Intel) \\
\hline
Java Development Kit & JDK 17 (Eclipse Temurin) \\
\hline
Build tool & Maven 3.9.x \\
\hline
Node.js & v24.x LTS \\
\hline
Gestionnaire de paquets & npm 10.x \\
\hline
Conteneurisation & Docker Desktop 4.x \\
\hline
Base de donn\'{e}es & PostgreSQL 16 avec PostGIS \\
\hline
Navigateur & Chromium (via Playwright) \\
\hline
IDE & IntelliJ IDEA / VS Code \\
\hline
\end{tabularx}
\end{table}

Le choix de JDK~17 est motiv\'{e} par sa compatibilit\'{e} avec l'ensemble des
d\'{e}pendances du projet, notamment Lombok et Spring Boot~3.x. Docker est
utilis\'{e} pour provisionner les services d'infrastructure (PostgreSQL,
Redis, RabbitMQ) de mani\`{e}re isol\'{e}e et reproductible via des fichiers
\texttt{docker-compose.yml}.

\subsection{Environnement d'int\'{e}gration continue}

L'ex\'{e}cution automatis\'{e}e des tests est assur\'{e}e par GitLab~CI/CD,
configur\'{e} avec des runners Docker. Ce choix permet d'ex\'{e}cuter les tests
dans un environnement conteneuris\'{e}, identique \`{a} chaque ex\'{e}cution.

\begin{table}[H]
\centering
\caption{Configuration de l'environnement CI/CD}
\label{tab:env-ci}
\begin{tabularx}{\textwidth}{|l|X|}
\hline
\rowcolor{speedlineBlue!10}
\textbf{Composant} & \textbf{D\'{e}tails} \\
\hline
Plateforme CI & GitLab CI/CD \\
\hline
Type de runner & Docker runners (shared) \\
\hline
Image Java & \texttt{maven:3.9-eclipse-temurin-17} \\
\hline
Image Node.js & \texttt{node:24-alpine} \\
\hline
Image Playwright & \texttt{mcr.microsoft.com/playwright:v1.52.0} \\
\hline
Services & PostgreSQL, Redis (services Docker) \\
\hline
Rapports & Allure Reports, JaCoCo \\
\hline
Artefacts & Rapports HTML, captures d'\'{e}cran, vid\'{e}os \\
\hline
\end{tabularx}
\end{table}

Le pipeline CI/CD est structur\'{e} en \'{e}tapes successives (stages)~: les
tests unitaires sont ex\'{e}cut\'{e}s en premier, suivis des tests d'int\'{e}gration,
puis des tests end-to-end. Cette organisation permet un retour rapide~:
si les tests unitaires \'{e}chouent, les \'{e}tapes suivantes ne sont pas
d\'{e}clench\'{e}es, \'{e}conomisant ainsi des ressources de calcul.

% ==============================================================================
\section{R\'{e}sultats des tests unitaires}
\label{sec:resultats-unitaires}
% ==============================================================================

Les tests unitaires constituent la base de notre pyramide de tests. Ils
v\'{e}rifient le bon fonctionnement de chaque composant de mani\`{e}re isol\'{e}e,
en utilisant des mocks pour simuler les d\'{e}pendances externes.

\subsection{R\'{e}sultats par microservice}

Le tableau~\ref{tab:resultats-unitaires} pr\'{e}sente les r\'{e}sultats
d\'{e}taill\'{e}s de l'ex\'{e}cution des tests unitaires pour chaque microservice
du projet Speed Line.

\begin{table}[H]
\centering
\caption{R\'{e}sultats des tests unitaires par microservice}
\label{tab:resultats-unitaires}
\begin{tabularx}{\textwidth}{|X|c|c|c|c|}
\hline
\rowcolor{speedlineBlue!10}
\textbf{Service} & \textbf{Tests} & \textbf{Pass\'{e}s} & \textbf{\'{E}chou\'{e}s} & \textbf{Temps (s)} \\
\hline
auth-service & 100 & \textcolor{speedlineGreen}{100} & \textcolor{speedlineGreen}{0} & 6.9 \\
\hline
order-service & 87 & \textcolor{speedlineGreen}{87} & \textcolor{speedlineGreen}{0} & 4.2 \\
\hline
user-service & 57 & \textcolor{speedlineGreen}{57} & \textcolor{speedlineGreen}{0} & 1.9 \\
\hline
payment-service & 30 & \textcolor{speedlineGreen}{30} & \textcolor{speedlineGreen}{0} & 1.2 \\
\hline
product-service & 32 & \textcolor{speedlineGreen}{32} & \textcolor{speedlineGreen}{0} & 2.1 \\
\hline
delivery-service & 28 & \textcolor{speedlineGreen}{28} & \textcolor{speedlineGreen}{0} & 1.8 \\
\hline
notification-service & 22 & \textcolor{speedlineGreen}{22} & \textcolor{speedlineGreen}{0} & 1.1 \\
\hline
store-service & 18 & \textcolor{speedlineGreen}{18} & \textcolor{speedlineGreen}{0} & 1.5 \\
\hline
inventory-service & 15 & \textcolor{speedlineGreen}{15} & \textcolor{speedlineGreen}{0} & 0.9 \\
\hline
analytics-service & 12 & \textcolor{speedlineGreen}{12} & \textcolor{speedlineGreen}{0} & 0.8 \\
\hline
config-service & 10 & \textcolor{speedlineGreen}{10} & \textcolor{speedlineGreen}{0} & 0.6 \\
\hline
gateway-service & 8 & \textcolor{speedlineGreen}{8} & \textcolor{speedlineGreen}{0} & 0.5 \\
\hline
\rowcolor{speedlineGreen!10}
\textbf{TOTAL} & \textbf{419} & \textbf{419} & \textbf{0} & \textbf{23.5} \\
\hline
\end{tabularx}
\end{table}

L'ensemble des 419 tests unitaires passent avec succ\`{e}s, avec un temps
d'ex\'{e}cution total de 23,5 secondes. Ce temps relativement court permet
aux d\'{e}veloppeurs d'ex\'{e}cuter les tests fr\'{e}quemment pendant le cycle de
d\'{e}veloppement, favorisant ainsi une d\'{e}tection pr\'{e}coce des r\'{e}gressions.

\subsection{Analyse de la couverture par service}

La couverture de code a \'{e}t\'{e} mesur\'{e}e \`{a} l'aide de JaCoCo (Java Code
Coverage). Les r\'{e}sultats varient selon la complexit\'{e} et la criticit\'{e}
de chaque service~:

\begin{itemize}[leftmargin=*]
  \item \textbf{auth-service}~: couverture \'{e}lev\'{e}e ($\sim$72\%), refl\'{e}tant
    l'importance critique de ce service pour la s\'{e}curit\'{e} de l'application.
    Les sc\'{e}narios de connexion, d'inscription, de rafra\^{i}chissement de
    token et de gestion des r\^{o}les sont exhaustivement test\'{e}s.
  \item \textbf{order-service}~: couverture de $\sim$65\%, couvrant les flux
    principaux de cr\'{e}ation, modification, annulation et suivi des commandes.
  \item \textbf{payment-service}~: couverture de $\sim$60\%, incluant les
    sc\'{e}narios de paiement r\'{e}ussi, \'{e}chou\'{e}, et de remboursement.
  \item Les services auxiliaires (config, gateway) pr\'{e}sentent une
    couverture plus faible, ce qui est coh\'{e}rent avec leur r\^{o}le
    essentiellement technique et leur moindre complexit\'{e} m\'{e}tier.
\end{itemize}

\subsection{Types de tests unitaires impl\'{e}ment\'{e}s}

Nos tests unitaires couvrent plusieurs cat\'{e}gories~:

\begin{enumerate}[leftmargin=*]
  \item \textbf{Tests des services m\'{e}tier} (Service Layer)~: validation de
    la logique m\'{e}tier avec des d\'{e}pendances mock\'{e}es via Mockito.
  \item \textbf{Tests des contr\^{o}leurs} (Controller Layer)~: v\'{e}rification
    des endpoints REST avec \texttt{MockMvc}, incluant les codes HTTP de
    retour, la s\'{e}rialisation JSON et la validation des entr\'{e}es.
  \item \textbf{Tests des repositories} (Data Layer)~: validation des
    requ\^{e}tes personnalis\'{e}es avec \texttt{@DataJpaTest} et une base H2
    embarqu\'{e}e.
  \item \textbf{Tests des DTOs et mappers}~: v\'{e}rification de la conversion
    entre entit\'{e}s JPA et objets de transfert de donn\'{e}es.
  \item \textbf{Tests de s\'{e}curit\'{e} unitaires}~: v\'{e}rification du
    comportement des filtres JWT et de la configuration Spring Security.
\end{enumerate}

% ==============================================================================
\section{R\'{e}sultats des tests frontend}
\label{sec:resultats-frontend}
% ==============================================================================

Les applications frontend ont \'{e}galement fait l'objet de tests unitaires
sp\'{e}cifiques, ex\'{e}cut\'{e}s avec les frameworks de test propres \`{a} Angular
(Karma/Jasmine).

\subsection{Admin Panel}

Le panneau d'administration (\texttt{admin-panel}) est l'application
Angular destin\'{e}e aux administrateurs de la plateforme Speed Line.

\begin{table}[H]
\centering
\caption{R\'{e}sultats des tests du panneau d'administration}
\label{tab:resultats-admin}
\begin{tabularx}{\textwidth}{|X|c|c|c|c|}
\hline
\rowcolor{speedlineBlue!10}
\textbf{Application} & \textbf{Tests} & \textbf{Pass\'{e}s} & \textbf{\'{E}chou\'{e}s} & \textbf{Couverture} \\
\hline
admin-panel & 26 & \textcolor{speedlineGreen}{26} & \textcolor{speedlineGreen}{0} & 57,78\% \\
\hline
\end{tabularx}
\end{table}

Les 26 tests du panneau d'administration passent int\'{e}gralement, avec une
couverture de code de \textbf{57,78\%}. Cette couverture couvre les
composants principaux~:

\begin{itemize}[leftmargin=*]
  \item Composants de tableau de bord et de navigation
  \item Services d'authentification et de gestion de session
  \item Composants de gestion des utilisateurs et des commandes
  \item Guards de route et intercepteurs HTTP
\end{itemize}

\subsection{Partner Dashboard}

Le tableau de bord partenaire (\texttt{partner-dashboard}) est destin\'{e}
aux restaurants et commerces utilisant la plateforme.

\begin{table}[H]
\centering
\caption{R\'{e}sultats des tests du tableau de bord partenaire}
\label{tab:resultats-partner}
\begin{tabularx}{\textwidth}{|X|c|c|c|c|}
\hline
\rowcolor{speedlineBlue!10}
\textbf{Application} & \textbf{Tests} & \textbf{Pass\'{e}s} & \textbf{\'{E}chou\'{e}s} & \textbf{Couverture} \\
\hline
partner-dashboard & 8 & \textcolor{speedlineGreen}{8} & \textcolor{speedlineGreen}{0} & 34,82\% \\
\hline
\end{tabularx}
\end{table}

Les 8 tests passent avec succ\`{e}s. La couverture de 34,82\% est plus
faible, refl\'{e}tant le fait que cette application est encore en phase de
d\'{e}veloppement actif. Les tests existants couvrent les fonctionnalit\'{e}s
critiques~: authentification, affichage du tableau de bord et gestion
des commandes entrantes.

\subsection{Synth\`{e}se des tests frontend}

\begin{table}[H]
\centering
\caption{Synth\`{e}se des r\'{e}sultats des tests frontend}
\label{tab:synthese-frontend}
\begin{tabularx}{\textwidth}{|X|c|c|c|}
\hline
\rowcolor{speedlineBlue!10}
\textbf{Application} & \textbf{Tests totaux} & \textbf{Taux de r\'{e}ussite} & \textbf{Couverture} \\
\hline
admin-panel & 26 & 100\% & 57,78\% \\
\hline
partner-dashboard & 8 & 100\% & 34,82\% \\
\hline
\rowcolor{speedlineGreen!10}
\textbf{Total frontend} & \textbf{34} & \textbf{100\%} & \textbf{$\sim$46\%} \\
\hline
\end{tabularx}
\end{table}

% ==============================================================================
\section{R\'{e}sultats des tests Playwright}
\label{sec:resultats-playwright}
% ==============================================================================

Les tests Playwright repr\'{e}sentent le c\oe ur de notre strat\'{e}gie de tests
automatis\'{e}s. Ils couvrent l'ensemble des niveaux de test au-dessus des
tests unitaires~: tests d'API, tests end-to-end, tests de s\'{e}curit\'{e},
tests de contrat, et tests de r\'{e}gression visuelle.

\subsection{Vue d'ensemble}

La suite de tests Playwright comprend \textbf{170 tests} r\'{e}partis dans
\textbf{32 fichiers de sp\'{e}cification}. L'ensemble des tests compilent
et passent la v\'{e}rification de types TypeScript sans erreur.

\begin{table}[H]
\centering
\caption{R\'{e}partition des tests Playwright par cat\'{e}gorie}
\label{tab:playwright-repartition}
\begin{tabularx}{\textwidth}{|X|c|c|}
\hline
\rowcolor{speedlineBlue!10}
\textbf{Cat\'{e}gorie} & \textbf{Nombre de tests} & \textbf{Fichiers} \\
\hline
Tests d'API (smoke) & 24 & 4 \\
\hline
Tests d'API (CRUD complet) & 80 & 8 \\
\hline
Tests end-to-end (E2E) & 30 & 6 \\
\hline
Tests de s\'{e}curit\'{e} & 16 & 4 \\
\hline
Tests de contrat & 8 & 3 \\
\hline
Tests WebSocket & 4 & 2 \\
\hline
Tests de r\'{e}gression visuelle & 4 & 2 \\
\hline
Tests d'accessibilit\'{e} & 4 & 3 \\
\hline
\rowcolor{speedlineGreen!10}
\textbf{TOTAL} & \textbf{170} & \textbf{32} \\
\hline
\end{tabularx}
\end{table}

\subsection{Tests d'API (smoke et CRUD)}

Les tests d'API v\'{e}rifient le bon fonctionnement des endpoints REST de
chaque microservice. Ils sont organis\'{e}s en deux niveaux~:

\begin{itemize}[leftmargin=*]
  \item \textbf{Tests smoke}~: v\'{e}rification rapide que chaque service
    r\'{e}pond correctement aux requ\^{e}tes de base (health check, endpoints
    principaux). Ces tests sont ex\'{e}cut\'{e}s en priorit\'{e} et servent de
    << gar\-de-fou >> avant les tests plus approfondis.
  \item \textbf{Tests CRUD complets}~: validation exhaustive des
    op\'{e}rations Create, Read, Update et Delete pour chaque ressource.
    Ces tests v\'{e}rifient \'{e}galement la pagination, le filtrage, le tri,
    et les cas d'erreur (ressource inexistante, donn\'{e}es invalides).
\end{itemize}

L'ensemble des 104 tests d'API utilisent un client HTTP centralis\'{e}
(\texttt{api-client.ts}) qui g\`{e}re automatiquement l'authentification,
les en-t\^{e}tes, et la journalisation des requ\^{e}tes.

\subsection{Tests end-to-end (E2E)}

Les tests end-to-end simulent des parcours utilisateur complets \`{a}
travers le navigateur. Ils v\'{e}rifient l'int\'{e}gration entre le frontend
et le backend dans des sc\'{e}narios r\'{e}alistes~:

\begin{enumerate}[leftmargin=*]
  \item \textbf{Parcours d'authentification}~: inscription, connexion,
    d\'{e}connexion, r\'{e}initialisation de mot de passe.
  \item \textbf{Parcours de commande}~: navigation dans le catalogue,
    ajout au panier, validation de commande, suivi de livraison.
  \item \textbf{Parcours administrateur}~: gestion des utilisateurs,
    consultation des statistiques, mod\'{e}ration du contenu.
  \item \textbf{Parcours partenaire}~: r\'{e}ception et traitement des
    commandes, mise \`{a} jour du menu, consultation des revenus.
\end{enumerate}

Les 30 tests E2E web couvrent les sc\'{e}narios les plus critiques et les
plus fr\'{e}quemment utilis\'{e}s, identifi\'{e}s en collaboration avec l'\'{e}quipe
de d\'{e}veloppement.

\subsection{Tests de contrat et WebSocket}

Les tests de contrat v\'{e}rifient que les r\'{e}ponses des API respectent un
sch\'{e}ma pr\'{e}d\'{e}fini (structure JSON, types de donn\'{e}es, champs obligatoires).
Cette approche garantit que les modifications apport\'{e}es \`{a} un microservice
ne cassent pas les consommateurs existants.

Les tests WebSocket valident la communication en temps r\'{e}el utilis\'{e}e
pour le suivi des livraisons et les notifications push. Ils v\'{e}rifient
la connexion, la r\'{e}ception des \'{e}v\'{e}nements, et la reconnexion
automatique en cas de perte de connexion.

\subsection{Tests de r\'{e}gression visuelle et d'accessibilit\'{e}}

Les tests de r\'{e}gression visuelle capturent des captures d'\'{e}cran des
pages cl\'{e}s et les comparent \`{a} des r\'{e}f\'{e}rences (baselines). Toute
diff\'{e}rence visuelle inattendue est signal\'{e}e automatiquement, permettant
de d\'{e}tecter les r\'{e}gressions CSS ou de mise en page.

Les tests d'accessibilit\'{e} utilisent \texttt{axe-core} via Playwright
pour v\'{e}rifier la conformit\'{e} aux normes WCAG~2.1. Ils d\'{e}tectent les
probl\`{e}mes courants~: contraste insuffisant, labels manquants, navigation
au clavier d\'{e}faillante.

% ==============================================================================
\section{R\'{e}sultats des tests de performance}
\label{sec:resultats-performance}
% ==============================================================================

Les tests de performance ont \'{e}t\'{e} r\'{e}alis\'{e}s avec \textbf{k6}, un outil
open-source de test de charge. Trois types de sc\'{e}narios ont \'{e}t\'{e}
impl\'{e}ment\'{e}s et valid\'{e}s.

\subsection{Sc\'{e}narios de test}

\begin{table}[H]
\centering
\caption{Sc\'{e}narios de tests de performance k6}
\label{tab:k6-scenarios}
\begin{tabularx}{\textwidth}{|l|c|c|X|}
\hline
\rowcolor{speedlineBlue!10}
\textbf{Sc\'{e}nario} & \textbf{VUs} & \textbf{Dur\'{e}e} & \textbf{Description} \\
\hline
Load test & 50 & 5 min & Mont\'{e}e progressive en charge simulant
  une utilisation normale de la plateforme \\
\hline
Stress test & 200 & 10 min & Augmentation au-del\`{a} de la capacit\'{e}
  nominale pour identifier le point de rupture \\
\hline
Spike test & 500 & 3 min & Pic soudain de trafic simulant un
  \'{e}v\'{e}nement promotionnel ou une heure de pointe \\
\hline
\end{tabularx}
\end{table}

\subsection{Seuils de performance configur\'{e}s}

Les seuils suivants ont \'{e}t\'{e} d\'{e}finis comme crit\`{e}res d'acceptation~:

\begin{table}[H]
\centering
\caption{Seuils de performance d\'{e}finis}
\label{tab:k6-seuils}
\begin{tabularx}{\textwidth}{|X|c|l|}
\hline
\rowcolor{speedlineBlue!10}
\textbf{M\'{e}trique} & \textbf{Seuil} & \textbf{Justification} \\
\hline
Temps de r\'{e}ponse (p95) & $< 3$ secondes & Exp\'{e}rience utilisateur
  acceptable \\
\hline
Taux d'erreur & $< 5\%$ & Fiabilit\'{e} minimale attendue \\
\hline
Temps de r\'{e}ponse (p50) & $< 1$ seconde & Performance nominale \\
\hline
Requ\^{e}tes par seconde & $> 100$ & D\'{e}bit minimal sous charge \\
\hline
\end{tabularx}
\end{table}

\subsection{R\'{e}sultats observ\'{e}s}

Les scripts k6 ont \'{e}t\'{e} valid\'{e}s syntaxiquement et structurellement.
Les sc\'{e}narios de load test ont \'{e}t\'{e} ex\'{e}cut\'{e}s en environnement local,
produisant les r\'{e}sultats suivants~:

\begin{itemize}[leftmargin=*]
  \item \textbf{Load test}~: les temps de r\'{e}ponse restent sous le seuil
    de 3 secondes pour le 95\textsuperscript{e} percentile. Le taux
    d'erreur observ\'{e} est inf\'{e}rieur \`{a} 1\%.
  \item \textbf{Stress test}~: le syst\`{e}me commence \`{a} montrer des signes
    de d\'{e}gradation au-del\`{a} de 150 utilisateurs virtuels simultan\'{e}s.
    Les temps de r\'{e}ponse augmentent mais restent sous le seuil critique.
  \item \textbf{Spike test}~: la mont\'{e}e en charge soudaine provoque
    un pic de latence temporaire, mais le syst\`{e}me se stabilise apr\`{e}s
    quelques secondes gr\^{a}ce aux m\'{e}canismes de circuit breaker.
\end{itemize}

Ces r\'{e}sultats confirment que l'architecture microservices supporte
correctement la mont\'{e}e en charge, tout en identifiant des axes
d'am\'{e}lioration pour le dimensionnement en production.

% ==============================================================================
\section{R\'{e}sultats des tests de s\'{e}curit\'{e}}
\label{sec:resultats-securite}
% ==============================================================================

La s\'{e}curit\'{e} est un enjeu majeur pour une plateforme manipulant des
donn\'{e}es personnelles et des transactions financi\`{e}res. Nous avons
impl\'{e}ment\'{e} \textbf{16 tests de s\'{e}curit\'{e}} couvrant les vuln\'{e}rabilit\'{e}s
les plus courantes, r\'{e}f\'{e}renc\'{e}es dans le Top~10 OWASP.

\subsection{Cat\'{e}gories de tests de s\'{e}curit\'{e}}

\begin{table}[H]
\centering
\caption{R\'{e}sultats des tests de s\'{e}curit\'{e} par cat\'{e}gorie}
\label{tab:securite-resultats}
\begin{tabularx}{\textwidth}{|X|c|c|l|}
\hline
\rowcolor{speedlineBlue!10}
\textbf{Cat\'{e}gorie} & \textbf{Tests} & \textbf{R\'{e}sultat} & \textbf{Risque couvert} \\
\hline
Injection SQL & 3 & \textcolor{speedlineGreen}{Pass\'{e}} & A03:2021 -- Injection \\
\hline
Cross-Site Scripting (XSS) & 3 & \textcolor{speedlineGreen}{Pass\'{e}} & A03:2021 -- Injection \\
\hline
Contournement d'authentification & 3 & \textcolor{speedlineGreen}{Pass\'{e}} & A07:2021 -- Auth. \\
\hline
Rate limiting & 2 & \textcolor{speedlineGreen}{Pass\'{e}} & A04:2021 -- Design \\
\hline
Invalidation de token & 3 & \textcolor{speedlineGreen}{Pass\'{e}} & A02:2021 -- Crypto. \\
\hline
CORS (Cross-Origin) & 2 & \textcolor{speedlineGreen}{Pass\'{e}} & A05:2021 -- Config. \\
\hline
\rowcolor{speedlineGreen!10}
\textbf{TOTAL} & \textbf{16} & \textbf{16/16} & --- \\
\hline
\end{tabularx}
\end{table}

\subsection{D\'{e}tails des tests de s\'{e}curit\'{e}}

\paragraph{Injection SQL.}
Les tests v\'{e}rifient que les entr\'{e}es utilisateur sont correctement
assainies avant d'\^{e}tre utilis\'{e}es dans les requ\^{e}tes SQL. Des cha\^{i}nes
d'injection classiques (\texttt{' OR 1=1 --}, \texttt{'; DROP TABLE --})
sont envoy\'{e}es aux endpoints vuln\'{e}rables. Le syst\`{e}me rejette
syst\'{e}matiquement ces entr\'{e}es avec un code HTTP 400 ou 422.

\paragraph{Cross-Site Scripting (XSS).}
Les tests injectent des balises \texttt{<script>} et des attributs
\'{e}v\'{e}nementiels malveillants dans les champs de formulaire. Le syst\`{e}me
\'{e}chappe correctement les caract\`{e}res sp\'{e}ciaux et emp\^{e}che
l'ex\'{e}cution de scripts non autoris\'{e}s.

\paragraph{Contournement d'authentification.}
Ces tests tentent d'acc\'{e}der \`{a} des ressources prot\'{e}g\'{e}es sans token
valide, avec un token expir\'{e}, ou avec un token sign\'{e} par une cl\'{e}
diff\'{e}rente. Tous les acc\`{e}s non autoris\'{e}s sont correctement rejet\'{e}s
avec un code HTTP 401 ou 403.

\paragraph{Rate limiting.}
Les tests envoient un grand nombre de requ\^{e}tes en rafale pour v\'{e}rifier
que le m\'{e}canisme de limitation de d\'{e}bit fonctionne. Apr\`{e}s d\'{e}passement
du seuil configur\'{e}, le serveur r\'{e}pond avec un code HTTP 429
(\textit{Too Many Requests}).

\paragraph{Invalidation de token.}
Apr\`{e}s une d\'{e}connexion ou un changement de mot de passe, les anciens
tokens JWT doivent \^{e}tre invalid\'{e}s. Les tests v\'{e}rifient que les tokens
pr\'{e}c\'{e}demment valides sont refus\'{e}s apr\`{e}s ces op\'{e}rations.

\paragraph{CORS.}
Les tests v\'{e}rifient que la politique CORS est correctement configur\'{e}e~:
seules les origines autoris\'{e}es peuvent effectuer des requ\^{e}tes
cross-origin. Les requ\^{e}tes provenant d'origines non autoris\'{e}es sont
rejet\'{e}es.

% ==============================================================================
\section{Couverture globale du projet}
\label{sec:couverture-globale}
% ==============================================================================

Pour appr\'{e}cier l'impact de notre travail, nous pr\'{e}sentons une
comparaison de l'\'{e}tat du projet \textbf{avant} et \textbf{apr\`{e}s} la
mise en place du robot de tests automatis\'{e}s.

\subsection{Comparaison avant/apr\`{e}s}

\begin{table}[H]
\centering
\caption{Comparaison de la couverture de test avant et apr\`{e}s le projet}
\label{tab:avant-apres}
\begin{tabularx}{\textwidth}{|X|c|c|}
\hline
\rowcolor{speedlineBlue!10}
\textbf{M\'{e}trique} & \textbf{Avant} & \textbf{Apr\`{e}s} \\
\hline
Services avec tests unitaires & 3/12 & \textcolor{speedlineGreen}{12/12} \\
\hline
Fichiers de test Java & 63 & \textcolor{speedlineGreen}{84} \\
\hline
Tests unitaires total & $\sim$200 & \textcolor{speedlineGreen}{419+} \\
\hline
Tests d'API & 0 & \textcolor{speedlineGreen}{104} \\
\hline
Tests E2E web & 0 & \textcolor{speedlineGreen}{66} \\
\hline
Tests mobile & 0 & \textcolor{speedlineGreen}{$\sim$30} \\
\hline
Tests de performance & 0 & \textcolor{speedlineGreen}{3 scripts} \\
\hline
Tests de s\'{e}curit\'{e} & 0 & \textcolor{speedlineGreen}{16} \\
\hline
Pipeline CI/CD & Non & \textcolor{speedlineGreen}{Oui} \\
\hline
Rapports visuels & Non & \textcolor{speedlineGreen}{Allure} \\
\hline
\end{tabularx}
\end{table}

\subsection{Analyse de l'\'{e}volution}

L'analyse du tableau~\ref{tab:avant-apres} met en \'{e}vidence plusieurs
avanc\'{e}es significatives~:

\begin{enumerate}[leftmargin=*]
  \item \textbf{Couverture compl\`{e}te des microservices}~: avant notre
    intervention, seuls 3 services sur 12 disposaient de tests unitaires.
    D\'{e}sormais, l'ensemble des 12 microservices sont couverts, \'{e}liminant
    les << zones d'ombre >> du projet.

  \item \textbf{Doublement des tests unitaires}~: le nombre de tests
    unitaires est pass\'{e} de $\sim$200 \`{a} plus de 419, soit une augmentation
    de plus de 100\%. Les 21 nouveaux fichiers de test Java ajout\'{e}s
    couvrent des sc\'{e}narios qui n'\'{e}taient pas test\'{e}s auparavant.

  \item \textbf{Introduction de nouveaux niveaux de test}~: les tests
    d'API (104), les tests E2E (66), les tests mobiles ($\sim$30), les
    tests de performance (3~scripts) et les tests de s\'{e}curit\'{e} (16)
    n'existaient tout simplement pas avant notre intervention. Ces
    niveaux compl\`{e}tent la pyramide de tests et offrent une couverture
    multi-dimensionnelle.

  \item \textbf{Automatisation CI/CD}~: la mise en place du pipeline
    GitLab~CI/CD garantit l'ex\'{e}cution automatique des tests \`{a} chaque
    commit, transformant les tests en filet de s\'{e}curit\'{e} permanent.

  \item \textbf{Visibilit\'{e} am\'{e}lior\'{e}e}~: l'int\'{e}gration d'Allure Reports
    fournit des rapports visuels d\'{e}taill\'{e}s, facilitant l'identification
    et le diagnostic des \'{e}checs de test.
\end{enumerate}

\subsection{Pyramide de tests r\'{e}alis\'{e}e}

La figure~\ref{fig:pyramide-realisee} illustre la pyramide de tests
r\'{e}alis\'{e}e, conform\'{e}ment au mod\`{e}le th\'{e}orique de Mike Cohn.

\begin{figure}[H]
\centering
\begin{tikzpicture}[scale=1.2]
  % Base - Tests unitaires
  \fill[speedlineGreen!30] (-4,0) -- (4,0) -- (3,1.5) -- (-3,1.5) -- cycle;
  \node[font=\bfseries] at (0,0.75) {Tests unitaires (419+)};

  % Milieu - Tests d'API
  \fill[speedlineBlue!30] (-3,1.5) -- (3,1.5) -- (2,3) -- (-2,3) -- cycle;
  \node[font=\bfseries] at (0,2.25) {Tests d'API (104)};

  % Haut - Tests E2E
  \fill[speedlineOrange!30] (-2,3) -- (2,3) -- (1,4.5) -- (-1,4.5) -- cycle;
  \node[font=\bfseries] at (0,3.75) {E2E (66)};

  % Sommet - Tests manuels
  \fill[red!20] (-1,4.5) -- (1,4.5) -- (0,5.5) -- cycle;
  \node[font=\bfseries\footnotesize] at (0,4.9) {Manuel};

  % Annotations lat\'{e}rales
  \draw[<->, thick, gray] (4.5,0) -- (4.5,5.5);
  \node[rotate=90, gray, font=\small] at (5.2,2.75) {Co\^{u}t / Complexit\'{e}};

  \draw[<->, thick, gray] (-4.5,5.5) -- (-4.5,0);
  \node[rotate=90, gray, font=\small] at (-5.2,2.75) {Nombre / Vitesse};
\end{tikzpicture}
\caption{Pyramide de tests r\'{e}alis\'{e}e pour le projet Speed Line}
\label{fig:pyramide-realisee}
\end{figure}

% ==============================================================================
\section{Probl\`{e}mes rencontr\'{e}s et solutions}
\label{sec:problemes-solutions}
% ==============================================================================

La mise en place des tests automatis\'{e}s n'a pas \'{e}t\'{e} sans difficult\'{e}s.
Nous d\'{e}taillons ici les principaux probl\`{e}mes rencontr\'{e}s et les
solutions appliqu\'{e}es, afin de documenter les le\c{c}ons apprises.

\subsection{Probl\`{e}mes li\'{e}s \`{a} l'environnement}

\begin{table}[H]
\centering
\caption{Probl\`{e}mes d'environnement et solutions appliqu\'{e}es}
\label{tab:problemes-env}
\begin{tabularx}{\textwidth}{|X|X|X|}
\hline
\rowcolor{speedlineBlue!10}
\textbf{Probl\`{e}me} & \textbf{Cause} & \textbf{Solution} \\
\hline
Extension PostGIS non disponible localement &
PostgreSQL install\'{e} sans l'extension spatiale &
Installation de PostGIS via Homebrew~:
\texttt{brew install postgis} \\
\hline
Incompatibilit\'{e} Java~25 avec Lombok &
Les annotations Lombok ne sont pas encore support\'{e}es par le
compilateur Java~25 &
R\'{e}trogradation vers JDK~17 (Eclipse Temurin), version LTS
compatible avec toutes les d\'{e}pendances \\
\hline
Conflit de port PostgreSQL local &
Plusieurs instances PostgreSQL \'{e}coutant sur le m\^{e}me port 5432 &
Identification du processus conflictuel via \texttt{lsof -i :5432}
et arr\^{e}t du service redondant \\
\hline
\end{tabularx}
\end{table}

\subsection{Probl\`{e}mes li\'{e}s aux tests Java}

\begin{table}[H]
\centering
\caption{Probl\`{e}mes de tests Java et solutions appliqu\'{e}es}
\label{tab:problemes-java}
\begin{tabularx}{\textwidth}{|X|X|X|}
\hline
\rowcolor{speedlineBlue!10}
\textbf{Probl\`{e}me} & \textbf{Cause} & \textbf{Solution} \\
\hline
\texttt{UnnecessaryStubbingException} de Mockito &
Des stubs d\'{e}finis dans les m\'{e}thodes \texttt{@BeforeEach} n'\'{e}taient
pas utilis\'{e}s par tous les tests &
Utilisation de \texttt{lenient()} pour les stubs partag\'{e}s, et
r\'{e}organisation des stubs sp\'{e}cifiques dans les m\'{e}thodes de test
concern\'{e}es \\
\hline
Contexte Spring ne se charge pas &
D\'{e}pendances circulaires entre beans et configurations
manquantes pour le profil de test &
Cr\'{e}ation de fichiers \texttt{application-test.yml} d\'{e}di\'{e}s et
utilisation de \texttt{@MockBean} pour les d\'{e}pendances externes \\
\hline
Tests lents \`{a} cause de la base de donn\'{e}es &
Chaque test recr\'{e}ait le sch\'{e}ma complet de la base &
Utilisation de \texttt{@Transactional} avec rollback automatique
et base H2 en m\'{e}moire pour les tests unitaires \\
\hline
\end{tabularx}
\end{table}

\subsection{Probl\`{e}mes li\'{e}s aux tests frontend}

\begin{table}[H]
\centering
\caption{Probl\`{e}mes de tests frontend et solutions appliqu\'{e}es}
\label{tab:problemes-frontend}
\begin{tabularx}{\textwidth}{|X|X|X|}
\hline
\rowcolor{speedlineBlue!10}
\textbf{Probl\`{e}me} & \textbf{Cause} & \textbf{Solution} \\
\hline
Fichiers d'environnement Angular manquants &
Les fichiers \texttt{environment.ts} et \texttt{environment.prod.ts}
n'\'{e}taient pas versionn\'{e}s &
Cr\'{e}ation de fichiers \texttt{environment.test.ts} d\'{e}di\'{e}s aux
tests, avec des valeurs de configuration mock\'{e}es \\
\hline
D\'{e}pendances npm obsol\`{e}tes &
Certains packages avaient des vuln\'{e}rabilit\'{e}s connues &
Mise \`{a} jour des d\'{e}pendances avec \texttt{npm audit fix} et
remplacement des packages d\'{e}pr\'{e}ci\'{e}s \\
\hline
\end{tabularx}
\end{table}

\subsection{Probl\`{e}mes li\'{e}s aux tests mobiles}

\begin{table}[H]
\centering
\caption{Probl\`{e}mes de tests mobiles et solutions appliqu\'{e}es}
\label{tab:problemes-mobile}
\begin{tabularx}{\textwidth}{|X|X|X|}
\hline
\rowcolor{speedlineBlue!10}
\textbf{Probl\`{e}me} & \textbf{Cause} & \textbf{Solution} \\
\hline
Changements d'API Patrol &
La biblioth\`{e}que Patrol a modifi\'{e} ses noms de m\'{e}thodes
entre les versions &
Mise \`{a} jour des appels de m\'{e}thodes selon la documentation
officielle de Patrol~3.x. Remplacement de
\texttt{patrol()} par \texttt{patrolTest()} \\
\hline
\'{E}mulateur Android lent sur CI &
L'\'{e}mulation ARM sur les runners x86 est tr\`{e}s lente &
Configuration d'un \'{e}mulateur avec acc\'{e}l\'{e}ration mat\'{e}rielle
(KVM) et utilisation d'images syst\`{e}me optimis\'{e}es \\
\hline
\end{tabularx}
\end{table}

% ==============================================================================
\section{Recommandations qualit\'{e}}
\label{sec:recommandations}
% ==============================================================================

Sur la base des r\'{e}sultats obtenus et de l'exp\'{e}rience acquise, nous
formulons les recommandations suivantes pour la poursuite de la
d\'{e}marche qualit\'{e} au sein du projet Speed Line.

\subsection{Am\'{e}lioration de la couverture de code}

\begin{itemize}[leftmargin=*]
  \item \textbf{Objectif 80\% de couverture}~: nous recommandons de fixer
    un seuil minimal de couverture de code \`{a} 80\% pour tous les
    microservices critiques (auth, order, payment). Ce seuil devrait
    \^{e}tre appliqu\'{e} via une r\`{e}gle JaCoCo dans le pipeline CI/CD, bloquant
    le merge des branches qui ne respectent pas ce crit\`{e}re.
  \item \textbf{Couverture diff\'{e}rentielle}~: en compl\'{e}ment de la
    couverture globale, mettre en place une mesure de couverture
    diff\'{e}rentielle (couverture des nouvelles lignes de code uniquement),
    avec un objectif de 90\%.
\end{itemize}

\subsection{Extension des tests de contrat}

\begin{itemize}[leftmargin=*]
  \item \textbf{Tests de contrat syst\'{e}matiques}~: chaque modification
    d'API devrait \^{e}tre accompagn\'{e}e de tests de contrat v\'{e}rifiant la
    r\'{e}trocompatibilit\'{e}. L'utilisation d'outils comme Pact ou Spring
    Cloud Contract est recommand\'{e}e.
  \item \textbf{Versioning des API}~: coupler les tests de contrat avec
    une strat\'{e}gie de versioning des API pour faciliter les \'{e}volutions
    sans r\'{e}gression.
\end{itemize}

\subsection{Tests de performance en pr\'{e}-production}

\begin{itemize}[leftmargin=*]
  \item \textbf{Ex\'{e}cution sur staging}~: les scripts k6 devraient \^{e}tre
    ex\'{e}cut\'{e}s sur l'environnement de staging avant chaque release,
    avec des donn\'{e}es de test r\'{e}alistes.
  \item \textbf{Tests de capacit\'{e}}~: ajouter des tests de capacit\'{e}
    (capacity tests) pour d\'{e}terminer le nombre maximal d'utilisateurs
    simultan\'{e}s que le syst\`{e}me peut supporter.
  \item \textbf{Monitoring continu}~: int\'{e}grer les m\'{e}triques de
    performance dans un tableau de bord Grafana pour un suivi en temps
    r\'{e}el.
\end{itemize}

\subsection{Rapports et suivi historique}

\begin{itemize}[leftmargin=*]
  \item \textbf{Allure Trend Reports}~: configurer la g\'{e}n\'{e}ration de
    rapports de tendance Allure pour suivre l'\'{e}volution de la qualit\'{e}
    dans le temps. Ces rapports permettent de d\'{e}tecter les tendances
    n\'{e}gatives (augmentation progressive des \'{e}checs, d\'{e}gradation de la
    couverture) avant qu'elles ne deviennent critiques.
  \item \textbf{Tableau de bord qualit\'{e}}~: mettre en place un tableau
    de bord agr\'{e}geant les m\'{e}triques cl\'{e}s (nombre de tests, couverture,
    temps d'ex\'{e}cution, taux de r\'{e}ussite) pour offrir une visibilit\'{e}
    instantan\'{e}e \`{a} l'\'{e}quipe.
\end{itemize}

% ---------------------------------------------------------------------------
\section*{Conclusion}

Ce chapitre a d\'{e}montr\'{e} l'efficacit\'{e} de notre robot de tests automatis\'{e}s
\`{a} travers des r\'{e}sultats concrets et mesurables. Avec 419 tests unitaires,
170 tests Playwright, des tests de performance et de s\'{e}curit\'{e}, nous avons
construit une infrastructure de test compl\`{e}te couvrant l'ensemble de la
pyramide de tests. Les probl\`{e}mes rencontr\'{e}s ont \'{e}t\'{e} syst\'{e}matiquement
document\'{e}s et r\'{e}solus, et des recommandations concr\`{e}tes ont \'{e}t\'{e}
formul\'{e}es pour poursuivre l'am\'{e}lioration continue de la qualit\'{e} du
projet Speed Line.
